% Figure: the Rankine half-body formed by a point source in a uniform stream,
% with the stagnation point, the dividing streamline, and the asymptotic
% half-width 2*pi*b = m/U.
\begin{figure}[htbp]
\centering
\begin{tikzpicture}
\begin{axis}[
  width=12cm, height=6.5cm,
  xlabel={$x/a$ (source at origin)},
  ylabel={$y/a$},
  xmin=-3, xmax=6, ymin=-3.0, ymax=3.0,
  axis lines=middle,
  axis line style={enggray},
  tick label style={font=\scriptsize},
  label style={font=\small},
  legend style={font=\scriptsize},
  legend pos=outer north east,
]
% dividing streamline psi = m/2 : r = a*(pi - theta)/sin(theta), here drawn
% as the closed-form body half-contour. Parametrise by theta in (0, pi).
% Upper half:
\addplot[engblue, very thick, samples=120, domain=0.06:3.08, variable=\t]
  ({ (3.14159 - \t)/sin(deg(\t)) * cos(deg(\t)) },
   { (3.14159 - \t)/sin(deg(\t)) * sin(deg(\t)) });
\addlegendentry{dividing streamline (body)}
% Lower half (mirror):
\addplot[engblue, very thick, samples=120, domain=0.06:3.08, variable=\t]
  ({ (3.14159 - \t)/sin(deg(\t)) * cos(deg(\t)) },
   { -(3.14159 - \t)/sin(deg(\t)) * sin(deg(\t)) });
% asymptotic half-width pi (in units of a): y -> +/- pi as x -> +inf
\addplot[engred, dashed, samples=2, domain=-3:6] {3.14159};
\addplot[engred, dashed, samples=2, domain=-3:6] {-3.14159};
\addlegendentry{asymptote $y=\pm\pi a$}
% freestream arrows
\draw[enggray, ->] (axis cs:-3,-2.4) -- (axis cs:-2.0,-2.4);
\draw[enggray, ->] (axis cs:-3,-1.2) -- (axis cs:-2.0,-1.2);
\draw[enggray, ->] (axis cs:-3,1.2) -- (axis cs:-2.0,1.2);
\draw[enggray, ->] (axis cs:-3,2.4) -- (axis cs:-2.0,2.4);
% stagnation point at x = -a
\filldraw[engblack] (axis cs:-1,0) circle (2.0pt);
\node[engblack, font=\scriptsize, anchor=south east] at (axis cs:-1,0.1) {stagnation $x=-a$};
% source marker
\node[englime, font=\scriptsize, anchor=south west] at (axis cs:0.05,0.1) {source};
\filldraw[englime] (axis cs:0,0) circle (1.8pt);
\end{axis}
\end{tikzpicture}
\caption[The Rankine half-body]{A point source of strength $m$ placed
in a uniform stream $U_\infty$ produces the open contour known as a
Rankine half-body. The dividing streamline (the contour the
freestream cannot cross) separates the source fluid inside from the
freestream fluid outside, so the contour can be read as the surface of
a solid nose of that shape. The flow stagnates a distance $a =
m/(2\pi U_\infty)$ upstream of the source, where the source's outward
push exactly cancels the freestream. Far downstream the body grows to
a constant half-width $\pi a$, because all the source volume flux must
eventually be carried inside the contour. The half-body is the
simplest model of a streamlined nose or a bridge pier facing a
current, and the surface pressure follows directly by Bernoulli once
the surface speed is known.}
\label{fig:vol03ch11:rankine-halfbody}
\end{figure}
